% si_figs.tex : generated by tools/make_si_figs.py (do not edit by hand).
% Supplementary Figs S7-S14; files outputs/figures/paper/v3_restructure/si/FigS<file>.pdf (S13 = file FigS14, S14 = file FigS15); captions from tools/si_fig_captions.py (trimmed legends).
% S7: FigS7.pdf (file time 2026-10-05 14:43); trimmed caption of FigS7_legend.md (180 words)
\begin{figure}[p]\centering
\includegraphics[width=\linewidth,height=0.72\textheight,keepaspectratio]{../../../outputs/figures/paper/v3_restructure/si/FigS7.pdf}
\caption{\textbf{Four scoring schemes in Alaska (M1 six-fold comparison): split maps, distances and RMSE.} This is the comparison cited in Supplementary Note 2, not the five-stage XH validation ladder of Fig.~1e. \textbf{a}--\textbf{d}, Fold assignment of the 13,606 labelled 1~km cells (six folds, seed 0) for random cells, 0.05° sites, 0.5° blocks and $k$-fold nearest-neighbour distance matching (kNNDM; 13 $k$-means clusters merged into six folds). \textbf{e}--\textbf{h}, Distances from test cells to the nearest training cell (blue) and from 20,000 sampled permafrost cells of the 0.02° prediction grid (mean annual air temperature below 0~°C) to the nearest label (grey), share of cells per logarithmic bin (seed 0); medians (mean of three seeds), 0.004, 1.11, 25.8 and 53.5~km test-to-training and 81.6~km map-to-label. \textbf{i}, RMSE (mean of three seeds) of direct ML (CatBoost), the Stefan model fitted to the training folds and the Stefan model plus a ridge residual (weight 0.75), rising from random cells to kNNDM from 11.6 to 17.9~cm, 14.2 to 14.5~cm and 12.77 to 13.70~cm. Alaska Albers equal-area projection; coastlines, Natural Earth 50~m.}\label{figsi:S7}
\end{figure}
\clearpage
% S8: FigS8.pdf (file time 2026-10-05 14:43); trimmed caption of FigS8_legend.md (182 words)
\begin{figure}[p]\centering
\includegraphics[width=\linewidth,height=0.72\textheight,keepaspectratio]{../../../outputs/figures/paper/v3_restructure/si/FigS8.pdf}
\caption{\textbf{Transfer without target labels: residual models by target and learner detail.} Error change relative to the source-coefficient Stefan model without target labels. \textbf{a}, \textbf{b}, Source anchor plus residual ML (CatBoost) with residual weight 0.25 (low-weight residual, \textbf{a}) and 1.0 (\textbf{b}) at each target's labelled-cell centre (17 targets; source excludes the parent region; post hoc re-analysis of the label-grid curves); below −40~cm, end colour (Tibetan Plateau in \textbf{b}, −47.34~cm); Alaska targets enlarged, Lena Delta targets offset with leader lines. \textbf{c}, Direct ML (ten learners, TabICL v2 with two context settings, CatBoost with physics inputs) and low-weight residual (three learners), four-region stratified mean and region rows, one axis per column; the main CatBoost uses the abstract-bundle contrast for the four regions and the label-grid curve for Alaska. Intervals, 95\% block-bootstrap confidence intervals (CIs) weighted by cell (thick) and by 0.5° block (thin); band, ±0.5~cm. Polar stereographic projections (central meridian 127°~E; Alaska enlargement, 157°~W), scale bars true at 70°~N; Tibetan Plateau inset, Lambert azimuthal equal-area. Permafrost zones, ESA CCI Permafrost fraction v4.0 (1997--2021 mean); coastlines, Natural Earth 50~m.}\label{figsi:S8}
\end{figure}
\clearpage
% S9: FigS9.pdf (file time 2026-10-05 14:43); trimmed caption of FigS9_legend.md (176 words)
\begin{figure}[p]\centering
\includegraphics[width=\linewidth,height=0.72\textheight,keepaspectratio]{../../../outputs/figures/paper/v3_restructure/si/FigS9.pdf}
\caption{\textbf{Coefficient recalibration: coefficient error and ten-label gain by target, soil-property Stefan curves and regional contrasts.} \textbf{a}, Source-coefficient error $\ln(E_{\mathrm{own}}/E_{\mathrm{source}})$, $E_{\mathrm{own}}$ fitted by least squares to all labels of the target (label-grid target table; positive, target coefficient larger; Tibetan Plateau not in the table, grey). \textbf{b}, Error change of the Stefan model recalibrated with ten target labels relative to the source coefficient (post hoc re-analysis of the label-grid curves); below −15~cm, end colour (Tibetan Plateau, −96.37~cm). \textbf{c}, Recalibrated and soil-property recalibrated Stefan models minus the source-coefficient Stefan model by number of labels, each region a new target (cell-weighted 95\% intervals). \textbf{d}, Four ten-label contrasts by region and four-region stratified mean (abstract contrast bundle): recalibrated minus source Stefan (−2.45~cm; 95\% CI −3.04 to −1.88), anchor plus residual ML minus recalibrated Stefan (−0.18; −0.53 to −0.01), pseudo-label augmentation of the anchor plus residual (−0.13; −0.31 to 0.22) and residual minus physics-input structure (−3.74; −5.23 to −2.75); 95\% block-bootstrap CIs weighted by cell (thick) and block (thin); band, ±0.5~cm. Map conventions as in Supplementary Fig.~S8.}\label{figsi:S9}
\end{figure}
\clearpage
% S10: FigS10.pdf (file time 2026-10-05 14:43); trimmed caption of FigS10_legend.md (179 words)
\begin{figure}[p]\centering
\includegraphics[width=\linewidth,height=0.72\textheight,keepaspectratio]{../../../outputs/figures/paper/v3_restructure/si/FigS10.pdf}
\caption{\textbf{Method selection: choices of the cross-validation rule by target and source cross-validation tuning of neural networks.} \textbf{a}, Method chosen most often by five-fold block cross-validation within ten target labels (25 selections per target and scoring mode, 5 splits × 5 draws; scoring modes pooled for sub-regions); the rule was designed after the label-grid results had been viewed. \textbf{b}, RMSE of the rule minus the fixed residual recipe (anchor plus residual ML, residual weight 0.25), ten labels, point estimates; beyond ±3~cm, end colours (Tibetan Plateau, −29.53~cm). \textbf{c}, Neural networks configured by leave-one-source-region-out cross-validation minus their default configuration, four-region stratified mean, for residual ML (weights 0.25 and 1.0) and direct ML, without labels and with ten and all labels; one axis per column. \textbf{d}, Change in source cross-validation error against change in target error without labels (32 points: four networks, two model kinds, four regions; Spearman $\rho = -0.32$, descriptive); open circles, the 12 points where tuning kept the default. Intervals, 95\% block-bootstrap CIs weighted by cell (thick) and block (thin); band, ±0.5~cm. Map conventions as in Supplementary Fig.~S8.}\label{figsi:S10}
\end{figure}
\clearpage
% S11: FigS11.pdf (file time 2026-10-05 14:44); trimmed caption of FigS11_legend.md (230 words)
\begin{figure}[p]\centering
\includegraphics[width=\linewidth,height=0.72\textheight,keepaspectratio]{../../../outputs/figures/paper/v3_restructure/si/FigS11.pdf}
\caption{\textbf{Gains within label-rich regions: error change by scoring block, high-resolution inputs, climate extrapolation and stacking.} All analyses were designed after earlier results had been viewed. \textbf{a}, \textbf{b}, Residual ML (cross-validated residual weight, mean of two seeds) minus the Stefan model recalibrated on the same labels, by 0.5° scoring block, all labels within regions (half of the blocks scored per split; Lena Delta 20 blocks over 24 splits, Canada 36 over 25; split means +0.47 and +0.45~cm); in \textbf{b} blocks are enlarged squares at their centres, and one Canadian block (+23.77~cm) takes the end colour. \textbf{c}, Residual ML with ten added within-grid inputs minus residual ML with the baseline inputs and minus recalibrated Stefan, three-region stratified mean (500 and 1000 labels: Alaska and Lena Delta only; auxiliary contrast, partly unblinded). \textbf{d}, Spatial proxy for climate extrapolation in Canada (warm blocks scored; not a true extrapolation test): penalty of direct ML relative to recalibrated Stefan, of residual ML relative to direct ML, and residual ML minus recalibrated Stefan in warm blocks; filled, main edition (extrapolation width 0.97 in $\sqrt{\mathrm{TDD}}$ units); open, basic edition (width −0.002). \textbf{e}, Label-weighted stacking of physics models and products minus recalibrated Stefan, and stacked residual minus residual, transfer pool and within regions (partly unblinded). Intervals, 95\% block-bootstrap CIs weighted by cell (thick) and block (thin); band, ±0.5~cm. Polar stereographic projections; coastlines, Natural Earth 50~m.}\label{figsi:S11}
\end{figure}
\clearpage
% S12: FigS12.pdf (file time 2026-10-05 14:47); trimmed caption of FigS12_legend.md (216 words)
\begin{figure}[p]\centering
\includegraphics[width=\linewidth,height=0.72\textheight,keepaspectratio]{../../../outputs/figures/paper/v3_restructure/si/FigS12.pdf}
\caption{\textbf{Placement of new labels: strategy contrasts by target, label proximity and learned placement policies.} \textbf{a}, \textbf{b}, Error change of residual ML (anchor plus residual, weight 0.25) with ten transfer labels placed by block stratification (\textbf{a}) or covariate maximin-distance spreading (\textbf{b}) instead of random cells (13 targets with a placement test; grey, not tested); above 4~cm, end colour (Tibetan Plateau in \textbf{b}, +12.71~cm). \textbf{c}, Proximity contrasts, four-region stratified mean (symbols, regions): direct ML on a random minus a block split; error relative to recalibrated Stefan with labels near the scored cells minus labels in held-out blocks, for direct and residual ML; residual ML with ten labels inside minus outside the scored blocks. \textbf{d}, Nearest minus farthest third of scored cells, for residual ML minus recalibrated Stefan and recalibrated minus source Stefan (cell-weighted 95\% intervals). \textbf{e}, Learned placement policy and three variants minus block stratification and minus covariate spreading (five-region mean, 10 labels; Lena Delta and Canada, 40 labels), learned policy minus random cells, and design-weighted coefficient minus learned policy in the Lena Delta (two-stage 95\% intervals; XD, exploratory). Panels \textbf{a}, \textbf{b} and \textbf{e} were designed after earlier results had been viewed. Intervals in \textbf{c} and \textbf{e}, 95\% CIs weighted by cell (thick) and block (thin); band, ±0.5~cm. Map conventions as in Supplementary Fig.~S8.}\label{figsi:S12}
\end{figure}
\clearpage
% S13: FigS14.pdf (file time 2026-10-05 14:44); trimmed caption of FigS14_legend.md (149 words)
\begin{figure}[p]\centering
\includegraphics[width=\linewidth,height=0.72\textheight,keepaspectratio]{../../../outputs/figures/paper/v3_restructure/si/FigS14.pdf}
\caption{\textbf{Prediction intervals for the Lena Delta without target labels.} \textbf{a}, ALT predicted by the source-coefficient Stefan model (source regions exclude the Lena Delta and a 100~km buffer; ERA5-Land 2015--2020 climatology). \textbf{b}, Width of the 90\% prediction interval of \textbf{a}: hierarchical conformal quantile of leave-one-region-out log errors ($q$ = 0.71; calibration regions Alaska, Canada, E Russia and W Russia), so the width is proportional to the prediction; with four calibration regions the interval carries no finite-sample guarantee. \textbf{c}, Number of covariates outside the 0.5--99.5 percentile range of the source rows, missing values included; all mapped cells have three or more (most often mean annual air temperature, freezing degree-days and coldest-month temperature); the count is not an error-risk indicator. 38,086 of 53,011 grid cells shown; 87 of 38,173 land cells (0.2\%) lack inputs. 1~km display grid; north polar stereographic projection centred at 126.7°~E; coastlines, Natural Earth 10~m.}\label{figsi:S13}
\end{figure}
\clearpage
% S14: FigS15.pdf (file time 2026-10-05 14:48); trimmed caption of FigS15_legend.md (214 words)
\begin{figure}[p]\centering
\includegraphics[width=\linewidth,height=0.72\textheight,keepaspectratio]{../../../outputs/figures/paper/v3_restructure/si/FigS15.pdf}
\caption{\textbf{Deployment scenarios: scenario matrix, sequential stopping rule and abstract contrast bundle.} \textbf{a}, Four-way verdict by stage (labels in the new region) and contrast for five regions and the four-region stratified mean (95\% block-bootstrap intervals weighted by cell and by block; equivalence margin ±0.5~cm); frames, non-inferior (margin 0.5~cm); hatching, more labels than the region provides; dagger, Lena Delta and Canada pool; source and recalibrated, source-coefficient and recalibrated Stefan models. \textbf{b}, Sequential stopping rule for recalibration (3, 10, then 40 labels; stop when the jackknife standard error of the log coefficient is at most $\tau$): RMSE change relative to always using 40 labels against the mean share of labels used, three-region mean (Lena Delta, Canada, Alaska; cell-weighted 95\% intervals) and regional values; with $\tau$ = 0.05 the rule used 0.91 of the labels and changed RMSE by −0.05~cm (95\% CI −0.08 to −0.002). \textbf{c}, Abstract contrast bundle in the order of Supplementary Table S14: four-region stratified mean (95\% block-bootstrap intervals by cell, thick, and block, thin), regional values (symbols) and wording after Holm correction over ten contrasts; ensemble, mean of scale-calibrated Stefan, Kudryavtsev and ESA CCI anchors; last row, interval scores in the Lena Delta, Canada and Alaska (+9.51~cm; 95\% CI −0.75 to 22.11). Horizontal axis linear within ±2~cm, logarithmic beyond.}\label{figsi:S14}
\end{figure}
\clearpage
